\documentclass[11pt,a4paper]{article}
\usepackage[T1]{fontenc}
\usepackage[utf8]{inputenc}
\usepackage{lmodern}
\usepackage{amsmath,amssymb,amsthm,mathtools}
\usepackage[margin=27mm,headheight=14pt]{geometry}
\usepackage{microtype}
\usepackage{booktabs,tabularx,array,longtable}
\usepackage{enumitem}
\usepackage{xcolor}
\usepackage{fancyhdr}
\usepackage{hyperref}
\usepackage{bookmark}
\usepackage[nameinlink,noabbrev]{cleveref}
\definecolor{linkblue}{RGB}{22,68,105}
\hypersetup{colorlinks=true,linkcolor=linkblue,citecolor=linkblue,urlcolor=linkblue,
  pdftitle={T228: A contractive nonlinear SGD model with a positive exact power tail},
  pdfauthor={T228 research campaign - synthesis with credited Math Brain proofs},
  pdfsubject={Concrete registry row T228; research proof with internal audit}}
\urlstyle{same}
\setlength{\parindent}{1.2em}
\setlength{\parskip}{0.2em}
\setlength{\emergencystretch}{1.5em}
\setlist[enumerate]{itemsep=0.35em,topsep=0.45em}
\setlist[itemize]{itemsep=0.2em,topsep=0.4em}
\numberwithin{equation}{section}
\newtheorem{theorem}{Theorem}[section]
\newtheorem{lemma}[theorem]{Lemma}
\newtheorem{proposition}[theorem]{Proposition}
\newtheorem{corollary}[theorem]{Corollary}
\theoremstyle{remark}
\newtheorem{remark}[theorem]{Remark}
\DeclareMathOperator{\Lip}{Lip}
\DeclareMathOperator{\Law}{Law}
\DeclareMathOperator{\Leb}{Leb}
\DeclareMathOperator{\sech}{sech}
\DeclareMathOperator{\supp}{supp}
\newcommand{\E}{\mathbb E}
\newcommand{\Pp}{\mathbb P}
\newcommand{\Q}{\mathbb Q}
\newcommand{\R}{\mathbb R}
\newcommand{\ind}{\mathbf 1}
\newcommand{\dd}{\,\mathrm d}
\newcommand{\abs}[1]{\lvert#1\rvert}
\newcommand{\bigabs}[1]{\left|#1\right|}
\newcommand{\sourcehash}[1]{{\small\texttt{#1}}}
\pagestyle{fancy}
\fancyhf{}
\fancyhead[L]{\small T228: nonlinear SGD with an exact power tail}
\fancyhead[R]{\small 23 October 2026}
\fancyfoot[C]{\thepage}
\renewcommand{\headrulewidth}{0.3pt}
\title{\vspace{-1.4em}\textbf{A contractive nonlinear SGD model\\with a positive exact power tail}\\[0.55em]
  \large A complete proof for the concrete-model requirement in registry row T228}
\author{\normalsize T228 research campaign --- synthesis with credited Math Brain proofs}
\date{23 October 2026}

\begin{document}
\maketitle
\thispagestyle{plain}
\begin{center}
\small\textbf{Research proof; independent internal audit, no external peer review}
\end{center}

\begin{abstract}
For the smooth nonlinear activation $\sigma(z)=z+10^{-3}\tanh z$, we study
ordinary constant-step, batch-one squared-loss SGD with Gaussian inputs,
independent Gaussian label noise, and regression teacher $\theta\in[-1,1]$.
For every fixed $\eta\in[0.99,1.01]$, the actual random update maps satisfy
$\E\Lip(\Psi_\theta)<0.99010707<1$. There is a unique stationary law,
geometric convergence in Wasserstein distance of order one, and a finite
strictly positive exact absolute-tail constant. The tail exponent is the
unique $\alpha\in(1,2)$ satisfying $\E|1-\eta A^2|^\alpha=1$; it does not
depend on the teacher. We verify every input to the scalar implicit renewal
theorem and give the full forward-excursion proof of strict positivity.
The theorem supplies a concrete nonlinear example for T228. It builds on
credited Math Brain proofs and makes no general deep-network theorem,
external acceptance, or novelty certification claim.
\end{abstract}

\section{Registry requirement, model, and main theorem}
\label{sec:main}

\paragraph{The T228 requirement (paraphrase).}
Choose a concrete nonlinear regression model with Gaussian inputs and
persistent label noise, run unmodified constant-step SGD at batch size $b$,
and prove existence of a stationary law, average contraction, asymptotic
linearity, a nonarithmetic multiplier, and a nonzero tail constant. Identify
the power-tail index through the multiplier's Lyapunov equation and exhibit
a nonempty parameter regime with index below two. This is an example
requirement, not a universal assertion about deep neural networks.

The registry is titled \emph{Heavy-tailed stationary behaviour of SGD on
neural-network losses}. Its motivation includes the linear-regression
results of G\"urb\"uzbalaban, \c{S}im\c{s}ekli, and Zhu
\cite{Gurbuzbalaban2021}. We use a smooth bounded \emph{perturbation of the
identity}; the activation itself is unbounded. The row's bounded-activation
network is an illustrative option, not the model proved here. Scope and
source attribution are recorded in \cref{sec:attribution}.

Fix, throughout a trajectory,
\begin{equation}
\delta=\frac1{1000},\qquad \sigma(z)=z+\delta\tanh z,\qquad
\theta\in[-1,1],\qquad \eta\in\left[\frac{99}{100},\frac{101}{100}\right],
\qquad b=1.\label{eq:model-parameters}
\end{equation}
Let $(A_k,E_k)_{k\geq1}$ be independent and identically distributed pairs,
each pair consisting of two independent standard normal variables. Set
$Y_k=\sigma(A_k\theta)+E_k$. For a fresh pair $(A,E)$ and its label
$Y=\sigma(A\theta)+E$, define
\begin{align}
 f_\theta(x;A,E)&=\frac12\bigl(\sigma(Ax)-\sigma(A\theta)-E\bigr)^2,
 \label{eq:loss}\\
 \Psi_\theta(x)&=x-\eta A\sigma'(Ax)\bigl(\sigma(Ax)-Y\bigr),
 \qquad X_{k+1}=\Psi_{\theta,k+1}(X_k).\label{eq:update}
\end{align}
This is ordinary, unclipped, constant-step SGD on the sample loss
\eqref{eq:loss}. For $\theta\ne0$, labels and inputs are dependent.

\begin{theorem}[Explicit nonlinear model for T228]\label{thm:main}
For every fixed pair $(\eta,\theta)$ in \eqref{eq:model-parameters}, the
following conclusions hold.
\begin{enumerate}[label=\textnormal{(\roman*)}]
\item The smooth sample maps are nonaffine almost surely, and
\begin{equation}
 \E\Lip(\Psi_\theta)<q:=\frac{99010707}{100000000}
 =0.99010707<1.\label{eq:q-main}
\end{equation}
They therefore contract on average, also in the logarithmic sense.
\item There is exactly one invariant probability law $\pi_{\eta,\theta}$,
including among laws without an assumed first moment. It has a finite
first moment, has support $\R$, and for every deterministic $x\in\R$,
\begin{equation}
 W_1\bigl(\Law(X_n^x),\pi_{\eta,\theta}\bigr)
 \le q^n W_1\bigl(\delta_x,\pi_{\eta,\theta}\bigr).
 \label{eq:wasserstein}
\end{equation}
Here $\delta_x$ is the unit point mass at $x$.
\item There is exactly one positive solution $\alpha=\alpha(\eta)$ of
\begin{equation}
 h_\eta(\alpha):=\E\abs{1-\eta A^2}^\alpha=1,
 \qquad 1<\alpha<2.\label{eq:alpha-main}
\end{equation}
The law of $\log|1-\eta A^2|$ is nonarithmetic. The exponent is independent
of $\theta$.
\item A finite strictly positive constant $C_{\eta,\theta}$ satisfies
\begin{equation}
 \lim_{t\to\infty}t^\alpha\pi_{\eta,\theta}
 \bigl(\{x:\abs{x}>t\}\bigr)=C_{\eta,\theta}\in(0,\infty).
 \label{eq:tail-main}
\end{equation}
The same limit and constant hold with $|x-\theta|$ in place of $|x|$.
For every $p>0$, both the absolute $p$th moment and the estimation-error
$p$th moment are finite exactly when $p<\alpha$. Their second moments
are infinite.
\end{enumerate}
If $X\sim\pi_{\eta,\theta}$ is independent of the fresh step $(A,E)$, then
\begin{equation}
 C_{\eta,\theta}
 =\frac{\E\!\left[\abs{\Psi_\theta(X)}^\alpha
                   -\abs{(1-\eta A^2)X}^\alpha\right]}
 {\alpha\,\E\!\left[\abs{1-\eta A^2}^\alpha
                         \log\abs{1-\eta A^2}\right]}.
 \label{eq:constant-main}
\end{equation}
The numerator is an absolutely integrable \emph{coupled difference}, not
a subtraction of two separately finite expectations.
\end{theorem}

\begin{remark}[Symmetry and uniformity]\label{rem:symmetry}
When $\theta=0$, the stationary law is symmetric and each one-sided tail
has constant $C_{\eta,0}/2$. No one-sided tail limit or equality of one-sided
constants is asserted for a fixed nonzero teacher. Indeed, its stationary
law is not symmetric about zero; see \cref{sec:symmetry}. The contraction
bound is uniform on the stated rectangle. No uniform positive lower bound
for $C_{\eta,\theta}$ is claimed.
\end{remark}

The proof of \cref{thm:main} occupies \cref{sec:bounds,sec:exponent,sec:invariance,sec:minorization,sec:renewal,sec:positive,sec:moments}.
The nonzero-signal example $\theta=\eta=1$ is made explicit in
\cref{sec:signal}.

\section{Bounds on the actual nonlinear update}\label{sec:bounds}

\begin{lemma}[Global derivative and uniform remainder]\label{lem:bounds}
Put $M=1-\eta A^2$ and $\xi=|M|$. Define
\begin{align}
 L_\theta&=\xi+\eta A^2(4\delta+3\delta^2+2\delta|Y|),
 \label{eq:L}\\
 B_\theta&=\eta|A|\bigl[\delta+(1+\delta)(\delta+|Y|)\bigr].
 \label{eq:B}
\end{align}
Then $\Lip(\Psi_\theta)\le L_\theta$ and
\begin{equation}
 \Psi_\theta(x)=Mx+R_\theta(x),\qquad
 \sup_{x\in\R}|R_\theta(x)|\le B_\theta.\label{eq:remainder}
\end{equation}
Both $L_\theta$ and $B_\theta$ have moments of every positive order. The
map $\Psi_\theta$ is nonaffine almost surely.
\end{lemma}

\begin{proof}
Write $z=Ax$, $t=\tanh z$, and $u=\sech^2 z$. Differentiating the
\emph{actual} update yields
\begin{align}
 \sigma'(z)&=1+\delta u,&\sigma''(z)&=-2\delta ut,\nonumber\\
 \Psi_\theta'(x)&=1-\eta A^2 H(z,Y),&
 H(z,Y)&=(1+\delta u)^2-2\delta ut(z+\delta t-Y).
 \label{eq:derivative}
\end{align}
We have $|t|\le1$, $0\le u\le1$, and $|z|u\le1$. For the last bound,
$\cosh^2 z=1+\sinh^2 z\ge1+z^2\ge|z|$. Expansion gives
\begin{align}
 |H-1|&\le2\delta u+\delta^2u^2+2\delta|z|u|t|
              +2\delta^2ut^2+2\delta u|t||Y|\nonumber\\
       &\le4\delta+3\delta^2+2\delta|Y|.\label{eq:H-bound}
\end{align}
The mean value theorem on $\R$ proves the global Lipschitz bound
\eqref{eq:L}. It bounds the actual nonlinear map, not merely its
asymptotic multiplier.

For the remainder, use the exact identity
\begin{equation}
 \sigma'(z)(\sigma(z)-Y)
   =z+\delta zu+(1+\delta u)(\delta t-Y).
 \label{eq:gradient-expansion}
\end{equation}
Multiplication by $\eta A$ and the preceding bounds give
\eqref{eq:remainder} with \eqref{eq:B}.

The label envelope keeps the input-label dependence:
\begin{equation}
 |Y|\le(1+\delta)|\theta||A|+|E|.
 \label{eq:label-envelope}
\end{equation}
In particular, throughout the parameter rectangle,
\begin{equation}
 B_\theta\le\frac{101}{100}|A|
 \bigl[\delta+(1+\delta)\bigl(\delta+(1+\delta)|A|+|E|\bigr)\bigr].
 \label{eq:uniform-B}
\end{equation}
The envelopes for $B_\theta$, $L_\theta$, and $Y$ are polynomials in
$|A|$ and $|E|$, hence have all positive moments. Neither $B_\theta$ nor
$Y$ is required to be independent of $M$.

At zero, the derivative is $1-\eta A^2(1+\delta)^2$. For $A\ne0$,
the terms involving $u$ or $zu$ vanish as $|x|\to\infty$, and
$\Psi_\theta'(x)\to1-\eta A^2$. These values differ for every $A\ne0$,
an event of probability one. Thus the map is nonaffine almost surely.
\end{proof}

\section{Exact contraction arithmetic and the exponent}\label{sec:exponent}

\begin{lemma}[Uniform average contraction]\label{lem:contraction}
The constant $q$ in \eqref{eq:q-main} bounds $\E L_\theta$ strictly. Also
$\E\log L_\theta$ exists as a finite number and is negative; in particular,
$\E\log\Lip(\Psi_\theta)<0$ in the extended sense.
\end{lemma}

\begin{proof}
Let $\varphi$ be the standard normal density. Splitting at $|A|=1$,
using $\E(1-A^2)=0$ and $(a\varphi(a))'=(1-a^2)\varphi(a)$, gives
\begin{equation}
 \E|1-A^2|=4\varphi(1)=\frac4{\sqrt{2\pi e}}<\frac{97}{100}.
 \label{eq:normal-identity}
\end{equation}
The strict inequality needs no numerical quadrature. The elementary bounds
$\pi>157/50$ and $e>1359/500$ imply
\begin{equation}
 \left(\frac{97}{100}\right)^2
 2\frac{157}{50}\frac{1359}{500}-16
 =\frac{7532467}{125000000}>0.\label{eq:rational-normal}
\end{equation}
Since $|1-\eta A^2|\le|1-A^2|+|\eta-1|A^2$, it follows that
\begin{equation}
 \E\xi<\frac{98}{100}.\label{eq:xi-first-moment}
\end{equation}

The mixed label moment must retain the dependence between $A$ and $Y$.
By \eqref{eq:label-envelope},
\begin{align}
 \E[A^2|Y|]
 &\le(1+\delta)|\theta|\E|A|^3+\E[A^2|E|]\nonumber\\
 &=(1+\delta)|\theta|\,2\sqrt{\frac2\pi}+\sqrt{\frac2\pi}
 \le(3+2\delta)\sqrt{\frac2\pi}
 <3+2\delta=\frac{1501}{500}.\label{eq:mixed-label}
\end{align}
Only the second summand uses independence, namely that of $A$ and $E$.
There is no factorization of $\E[A^2|Y|]$. Thus
\begin{align}
 \E\Lip(\Psi_\theta)\le\E L_\theta
 &<\frac{98}{100}+\frac{101}{100}
  \left[\frac4{1000}+\frac3{10^6}
             +\frac2{1000}\frac{1501}{500}\right]\nonumber\\
 &=\frac{98}{100}+\frac{101}{100}\frac{10007}{10^6}
 =\frac{99010707}{10^8}=q<1.\label{eq:q-proof}
\end{align}
The exact margin is $1-q=989293/10^8=0.00989293$.

Jensen's inequality gives
$\E\log\Lip(\Psi_\theta)\le\log\E\Lip(\Psi_\theta)<\log q<0$,
allowing a negative infinite expectation. For the envelope itself, set
$c_0=4\delta+3\delta^2<1$. With $v=\eta A^2\ge0$,
$L_\theta\ge|1-v|+c_0v\ge c_0$, whereas
$\log^+L_\theta\le L_\theta$. Hence $\E|\log L_\theta|<\infty$ and
$\E\log L_\theta<\log q<0$.
\end{proof}

\begin{lemma}[Lyapunov root and logarithmic coefficient]\label{lem:root}
There is a unique positive root $\alpha\in(1,2)$ of
$h_\eta(p)=\E\xi^p=1$. For $0<p<\alpha$, $h_\eta(p)<1$. Moreover
\begin{equation}
 m_\alpha:=\E[\xi^\alpha\log\xi]\in(0,\infty),
 \label{eq:m-alpha}
\end{equation}
and $\log\xi$ is nonarithmetic.
\end{lemma}

\begin{proof}
The distribution of $M=1-\eta A^2$ is absolutely continuous. Its zeros
as a function of $A$ are $\pm\eta^{-1/2}$, with nonzero derivatives
there. Local logarithmic integrability near these roots and Gaussian decay
at infinity give
\begin{equation}
 \Pp(\xi=0)=0,\qquad \E|\log\xi|<\infty,\qquad
 \E[\xi^p|\log\xi|]<\infty\quad(p>0).
 \label{eq:log-moments}
\end{equation}
The law of $\log\xi$ is absolutely continuous and thus nonarithmetic.
Also $\E\log\xi\le\log\E\xi<0$.

The function $h_\eta$ is finite for all $p\ge0$, and its logarithm is
strictly convex because $\log\xi$ is nonconstant. We have
\begin{align}
 h_\eta(0)&=1,& h_\eta(1)&<0.98<1,\nonumber\\
 h_\eta(2)&=1-2\eta+3\eta^2
 \ge1-2\frac{99}{100}+3\left(\frac{99}{100}\right)^2
 =\frac{19603}{10000}>1.\label{eq:second-multiplier}
\end{align}
The quadratic is increasing on the permitted interval. Since
$\Pp(\xi>1)>0$, $h_\eta(p)\to\infty$ as $p\to\infty$.
Strict log-convexity therefore yields exactly one positive root, lying
between one and two, and $h_\eta(p)<1$ for $0<p<\alpha$.
The integrability in \eqref{eq:log-moments} makes $m_\alpha$ finite.
The derivative of $\log h_\eta$ at $\alpha$ is at least its strictly
positive secant slope from $1$ to $\alpha$, so $m_\alpha>0$.
All of these quantities depend on $\eta$ and not on $\theta$.
\end{proof}

Equivalently, the equation specifying the index is
\begin{equation}
 \frac1{\sqrt{2\pi}}\int_\R |1-\eta a^2|^\alpha
            e^{-a^2/2}\dd a=1.\label{eq:alpha-integral}
\end{equation}
For a general batch $b$, the corresponding scalar asymptotic multiplier
is $1-(\eta/b)\chi_b^2$. The present theorem fixes $b=1$; the interval
in \eqref{eq:model-parameters} already supplies a nonempty regime.

\section{Invariant law, uniqueness, and lower moments}\label{sec:invariance}

\begin{proposition}[Invariant law and convergence]\label{prop:invariant}
There is a unique invariant probability $\pi=\pi_{\eta,\theta}$ on $\R$.
It has a finite first moment, satisfies \eqref{eq:wasserstein}, and
\begin{equation}
 \E_\pi|X|^p<\infty\qquad(0<p<\alpha).\label{eq:lower-moments}
\end{equation}
\end{proposition}

\begin{proof}
At zero,
\begin{equation}
 \Psi_\theta(0)=\eta(1+\delta)AY,\qquad
 \E|\Psi_\theta(0)|<\infty.\label{eq:update-zero}
\end{equation}
The Markov operator consequently maps probability laws with finite first
moment into themselves. Coupling two inputs by the same independent fresh
step gives
\begin{equation}
 W_1(\mu P_\theta,\nu P_\theta)
 \le(\E L_\theta)W_1(\mu,\nu)\le qW_1(\mu,\nu).
 \label{eq:W-contraction}
\end{equation}
Completeness of the space of such laws for $W_1$ and the contraction
mapping theorem provide an invariant law in that space and the geometric
bound \eqref{eq:wasserstein}.

To prove uniqueness without a hidden moment restriction, start a second
chain from any finite random value $Z$, independent of its future steps,
and couple it to the chain from zero using the same steps. Then
\begin{equation}
 |X_n^Z-X_n^0|\le |Z|\prod_{j=1}^n L_{\theta,j}
   \longrightarrow0\quad\text{almost surely}.\label{eq:path-coupling}
\end{equation}
Indeed, the product has expectation $(\E L_\theta)^n\le q^n$;
Markov's inequality followed by Borel--Cantelli proves that the product
tends to zero. If $Z$ has any invariant law, bounded Lipschitz test
functions and dominated convergence identify that law with $\pi$.
No moment of $Z$ is used.

For $0<p\le1$ with $p<\alpha$, \eqref{eq:remainder} and subadditivity
give
\begin{equation}
 \E|\Psi_\theta(x)|^p\le h_\eta(p)|x|^p+\E B_\theta^p.
 \label{eq:drift-small-p}
\end{equation}
For $1<p<\alpha$, use
$(a+b)^p\le(1+\varepsilon)a^p+C_{p,\varepsilon}b^p$ and choose
$\varepsilon>0$ so that $r_p=(1+\varepsilon)h_\eta(p)<1$.
In both cases this yields
\begin{equation}
 \E|\Psi_\theta(x)|^p\le r_p|x|^p+d_p,
 \qquad r_p<1,\quad d_p<\infty.\label{eq:moment-drift}
\end{equation}
The chain started from zero has uniformly bounded $p$th moments. Its
weak convergence to $\pi$ and the lower semicontinuity of $|x|^p$ imply
\eqref{eq:lower-moments}.
\end{proof}

\section{Explicit compact minorization and full support}\label{sec:minorization}

\begin{lemma}[Compact-to-compact minorization]\label{lem:minorization}
For every finite $R,T>0$ there is an explicit $\varepsilon_{R,T}>0$,
uniform over the parameter rectangle, such that for every Borel $F$,
\begin{equation}
 P_\theta(x,F)\ge\varepsilon_{R,T}
        \Leb(F\cap[-T,T])\qquad(|x|\le R).\label{eq:minorization}
\end{equation}
Consequently $\supp\pi_{\eta,\theta}=\R$.
\end{lemma}

\begin{proof}
Conditionally on $A=a$, the fresh $E$ remains standard normal and
\begin{align}
 \Psi_\theta(x)&=\mu_\theta(x,a)+\eta a\sigma'(ax)E,\nonumber\\
 \mu_\theta(x,a)&=x-\eta a\sigma'(ax)
                    [\sigma(ax)-\sigma(a\theta)].\label{eq:conditional-mean}
\end{align}
For $a\in[1,2]$ and $|x|\le R$, the inequalities
$\sigma'\le1+\delta$ and $|\sigma(v)|\le(1+\delta)|v|$ imply
\begin{equation}
 |\mu_\theta(x,a)|\le D_R
 :=R+4\frac{101}{100}(1+\delta)^2(R+1).\label{eq:D-R}
\end{equation}
The conditional standard deviation belongs to
\begin{equation}
 [s_{\min},s_{\max}]
 :=\left[\frac{99}{100},\;2\frac{101}{100}(1+\delta)\right].
 \label{eq:s-bounds}
\end{equation}
At each $|y|\le T$ its density is therefore at least
\begin{equation}
 d_{R,T}:=\frac1{s_{\max}\sqrt{2\pi}}
 \exp\!\left[-\frac{(T+D_R)^2}{2s_{\min}^2}\right]>0.
 \label{eq:density-bound}
\end{equation}
Integrate over $a\in[1,2]$ and set
$\varepsilon_{R,T}=\Pp(1\le A\le2)d_{R,T}>0$ to obtain
\eqref{eq:minorization}. In particular, every nonempty open interval
has positive transition probability from every point. Integrating this
strictly positive transition probability against $\pi$ proves that every
such interval has positive stationary mass.
\end{proof}

\section{Exact absolute-tail limit and the constant formula}\label{sec:renewal}

The external tail-limit input is the $N=1$ specialization of
Jelenkovi\'c--Olvera-Cravioto \cite[Theorem 3.1, printed p.~5]{JO2012Thm31}.
This is a discrepancy theorem for a nonnegative variable independent of
its coefficient. Its use does not require the nonlinear update itself
to be an affine fixed-point recursion.

\begin{lemma}[The scalar renewal input]\label{lem:renewal-input}
Let $R_0\ge0$ be independent of $\xi>0$. Suppose $\log\xi$ is
nonarithmetic, $\E\xi^\alpha=1$,
$0<\E[\xi^\alpha\log\xi]<\infty$, $\E\xi^\gamma<\infty$ for some
$0\le\gamma<\alpha$, and $\E R_0^p<\infty$ for every $0<p<\alpha$.
If
\begin{equation}
 \int_0^\infty t^{\alpha-1}
  \bigabs{\Pp(R_0>t)-\Pp(\xi R_0>t)}\dd t<\infty,
 \label{eq:renewal-condition}
\end{equation}
then $\lim_{t\to\infty}t^\alpha\Pp(R_0>t)=C\in[0,\infty)$, where
\begin{equation}
 C=\frac1{m_\alpha}\int_0^\infty t^{\alpha-1}
       [\Pp(R_0>t)-\Pp(\xi R_0>t)]\dd t.
 \label{eq:renewal-constant}
\end{equation}
\end{lemma}

\begin{proof}
Apply \cite[Theorem 3.1]{JO2012Thm31} with $N=1$ and $C_1=\xi$.
The stated conditions are its scalar hypotheses. The conclusion is
expressed as a scaled limit, also covering the possibility $C=0$.
\end{proof}

\begin{proposition}[Existence of the finite tail limit]\label{prop:tail-exists}
The limit in \eqref{eq:tail-main} exists in $[0,\infty)$ and has the
formula \eqref{eq:constant-main}. Its numerator is absolutely integrable.
\end{proposition}

\begin{proof}
Choose $X\sim\pi$ independently of the fresh $(A,E)$ and take
$R_0=|X|$. Then $R_0$ is independent of $\xi$. Lemma~\ref{lem:root}
provides the multiplier equation, logarithmic nonarithmeticity, and
$0<m_\alpha<\infty$. Taking $\gamma=0$ supplies the auxiliary coefficient
moment. Proposition~\ref{prop:invariant} supplies every stationary moment
below $\alpha$. Thus only \eqref{eq:renewal-condition} remains.

Let $U=|\Psi_\theta(X)|$ and $W=|MX|=\xi|X|$. Stationarity gives $U$
the law of $R_0$, and \eqref{eq:remainder} gives $|U-W|\le B_\theta$.
Tonelli's theorem and the indicator-tail integral yield
\begin{align}
 &\int_0^\infty t^{\alpha-1}
   \bigabs{\Pp(R_0>t)-\Pp(\xi R_0>t)}\dd t\nonumber\\
 &\hspace{2em}\le
 \E\int_0^\infty t^{\alpha-1}
       |\ind_{\{U>t\}}-\ind_{\{W>t\}}|\dd t
 =\frac1\alpha\E|U^\alpha-W^\alpha|.\label{eq:tail-discrepancy}
\end{align}
Because $1<\alpha<2$, the mean value theorem and subadditivity at
exponent $\alpha-1$ imply
\begin{align}
 |U^\alpha-W^\alpha|
 &\le\alpha B_\theta(W+B_\theta)^{\alpha-1}\nonumber\\
 &\le\alpha\left[B_\theta^\alpha
        +B_\theta\xi^{\alpha-1}|X|^{\alpha-1}\right].
 \label{eq:coupled-difference}
\end{align}
The right side is integrable. The stationary moment of order $\alpha-1$
is finite; $X$ is independent of the step; and the polynomial Gaussian
envelopes give $\E[B_\theta\xi^{\alpha-1}]<\infty$. This is a
\emph{joint} moment: independence of $B_\theta$ and $\xi$ is neither
asserted nor needed.

Lemma~\ref{lem:renewal-input} now gives a finite nonnegative limit.
Signed Fubini is legitimate by \eqref{eq:coupled-difference}, and gives
\begin{equation}
 \int_0^\infty t^{\alpha-1}
       [\Pp(U>t)-\Pp(W>t)]\dd t
 =\frac1\alpha\E[U^\alpha-W^\alpha].\label{eq:signed-fubini}
\end{equation}
Substitution into \eqref{eq:renewal-constant} proves
\eqref{eq:constant-main}. Strict positivity requires the separate
argument that follows.
\end{proof}

\section{Strict positivity by forward excursions}\label{sec:positive}

This section gives the complete positivity mechanism of the existing T228
work of \texttt{claude-b1e26540}, in particular
\texttt{T3\_STEP3.md} \cite{BrainT3Positive}, with every bound checked for
the nonzero-teacher law. The steps $(A_j,E_j)$ are iid, so
$(\xi_j,B_{\theta,j})$ are iid pairs. Their two components may be
dependent. The tilt below acts on the \emph{complete step}; no teacher
symmetry or rotation argument is used.

\subsection{Radial comparison and the stopped likelihood ratio}

Set
\begin{equation}
 S_0=0,\qquad S_n=\sum_{j=1}^n\log\xi_j,\qquad
 V_n=\sum_{j=1}^n e^{-S_j}B_{\theta,j}.\label{eq:S-V}
\end{equation}
Repeated reverse triangle inequalities in \eqref{eq:remainder} give,
pathwise,
\begin{equation}
 |X_n|\ge e^{S_n}(|X_0|-V_n).\label{eq:radial-lower}
\end{equation}
To see this directly, multiply
$|X_j|\ge\xi_j|X_{j-1}|-B_{\theta,j}$ by $e^{-S_j}$ and sum the
resulting telescoping inequalities.

Tilt each complete step by $\xi^\alpha$ and call the iid product law
$\Q$. It is a probability law because $\E\xi^\alpha=1$.
On the sigma-field $\mathcal F_n$ of the first $n$ complete steps,
$\mathrm d\Q/\mathrm d\Pp=e^{\alpha S_n}$. For any stopping time $\tau$
and stopped event $F\in\mathcal F_\tau$, summing over the disjoint
events $\{\tau=n\}$ proves
\begin{equation}
 \Pp(F,\tau<\infty)
 =\E_\Q\!\left[e^{-\alpha S_\tau}
                    \ind_{F\cap\{\tau<\infty\}}\right].
 \label{eq:stopped-change}
\end{equation}
This identity makes no infinite-path absolute-continuity or uniform-
integrability assertion. Under $\Q$, the increment $Z=\log\xi$ satisfies
$\E_\Q|Z|<\infty$ and $\E_\Q Z=m_\alpha>0$. Hence
$S_n\to\infty$ almost surely under $\Q$. Fix
\begin{equation}
 \beta=\frac12,\qquad
 r:=\E_\Q e^{-\beta Z}=h_\eta(\alpha-\beta)<1.
 \label{eq:r-tilt}
\end{equation}
The inequality holds because $0<\alpha-\beta<\alpha$.

\subsection{Positive survival, tight overshoots, and finite remainder}

\begin{lemma}[Three tilted-walk estimates]\label{lem:tilted-estimates}
Under $\Q$, the following hold:
\begin{enumerate}[label=\textnormal{(\alph*)}]
\item The positive-survival probability obeys
\begin{equation}
 p_+:=\Q(S_j>0\text{ for every }j\ge1)
 \ge p_0:=\Q(Z>1)(1-e^{-\beta})>0.\label{eq:survival}
\end{equation}
\item If $\tau_u=\inf\{n\ge1:S_n>u\}$ for $u\ge0$ and
$O_u=S_{\tau_u}-u$, then the overshoots $O_u$ are uniformly tight:
\begin{equation}
 \sup_{u\ge0}\Q(O_u>z)
 \le K_U\left[\Q(Z>z)+\int_z^\infty\Q(Z>v)\dd v\right]
 \longrightarrow0,\qquad z\to\infty,\label{eq:overshoot-final}
\end{equation}
where $K_U=e^\beta/(1-r)$.
\item The discounted remainder is finite almost surely, and
\begin{equation}
 \E_\Q V_\infty^\beta
 \le\frac{\E_\Pp[\xi^{\alpha-\beta}B_\theta^\beta]}
             {1-h_\eta(\alpha-\beta)}<\infty,
 \qquad V_\infty:=\sum_{j\ge1}e^{-S_j}B_{\theta,j}.
 \label{eq:V-infinite}
\end{equation}
\end{enumerate}
\end{lemma}

\begin{proof}
\emph{(a) Survival.} By \eqref{eq:r-tilt}, $e^{-\beta S_n}$ is a
nonnegative supermartingale. Stop it at the first time an independent
walk reaches $(-\infty,-1]$, truncated at $n$. On the event that this
time is at most $n$, its stopped value is at least $e^\beta$.
The supermartingale inequality and then monotone convergence give
\begin{equation}
 \Q\!\left(\inf_{n\ge0}S_n\le-1\right)\le e^{-\beta}.
 \label{eq:survival-ruin}
\end{equation}
A first increment exceeding one followed by a fresh walk that never
reaches $-1$ stays strictly positive. Independence proves
\eqref{eq:survival}. The Gaussian input is unbounded, and the tilt has
strictly positive density where $\xi>0$, so $\Q(Z>1)>0$.

\emph{(b) Overshoots.} Positive drift makes $\tau_u$ finite almost surely
under $\Q$. Chernoff's inequality and \eqref{eq:r-tilt} give
\begin{equation}
 \sum_{n\ge0}\Q(S_n\le1)
 \le\sum_{n\ge0}e^\beta r^n
 =\frac{e^\beta}{1-r}=K_U.\label{eq:occupation-base}
\end{equation}
For $I=(a,a+1]$, let $\nu_I$ be its first entry time, allowing time zero.
Conditionally on $S_{\nu_I}\in I$, a later visit to $I$ entails a
fresh-walk displacement at most one. The strong Markov property and
\eqref{eq:occupation-base} therefore give the translated occupation bound
\begin{equation}
 U(I):=\sum_{n\ge0}\Q(S_n\in I)\le K_U.\label{eq:occupation-translated}
\end{equation}
Partition $(-\infty,u]$ into the intervals $(u-k-1,u-k]$, $k\ge0$.
Dropping the restriction that no earlier crossing has occurred, and using
independence of the next increment from the past, yields for $z\ge0$
\begin{align}
 \Q(O_u>z)
 &\le\sum_{n\ge1}\E_\Q\!
  \left[\ind_{\{S_{n-1}\le u\}}\,
        \Q(Z>u+z-S_{n-1})\right]\nonumber\\
 &\le K_U\sum_{k\ge0}\Q(Z>z+k)\nonumber\\
 &\le K_U\left[\Q(Z>z)+\int_z^\infty\Q(Z>v)\dd v\right].
 \label{eq:overshoot-proof}
\end{align}
The last expression is independent of $u$ and tends to zero because
$\E_\Q Z^+<\infty$.

\emph{(c) Discounted remainders.} Subadditivity at $\beta=1/2$, Tonelli's
theorem, and independence of a complete step from its past imply
\begin{align}
 \E_\Q\left(\sum_{j\ge1}e^{-S_j}B_{\theta,j}\right)^\beta
 &\le\sum_{j\ge1}r^{j-1}\E_\Q[\xi^{-\beta}B_\theta^\beta]\nonumber\\
 &=\frac{\E_\Pp[\xi^{\alpha-\beta}B_\theta^\beta]}
              {1-h_\eta(\alpha-\beta)}<\infty.
 \label{eq:discounted-proof}
\end{align}
The numerator is the correct joint moment. It is finite by the polynomial
Gaussian envelopes and $\alpha-\beta>0$; there is no inverse-multiplier
singularity. Finite $\beta$th moment implies $V_\infty<\infty$
almost surely under $\Q$.
\end{proof}

\subsection{A uniform excursion lower bound}

\begin{lemma}[Annulus-to-tail excursion]\label{lem:excursion}
There exist $K\ge1$ and $z\ge0$, depending on the fixed parameters but
not on the tail threshold, such that for $t>4K$,
\begin{equation}
 \inf_{2K\le|y|\le4K}\Pp_y(\tau_{D_t}<\tau_C)
 \ge\frac{p_0}{2}e^{-\alpha z}\left(\frac Kt\right)^\alpha,
 \quad C=[-K,K],\quad D_t=\{x:|x|>t\}.
 \label{eq:annulus-bound}
\end{equation}
The hitting times here start at time one.
\end{lemma}

\begin{proof}
Choose $z\ge0$ so that \eqref{eq:overshoot-final} is at most $p_0/4$.
Next choose $K\ge1$ so that $\Q(V_\infty>K/2)\le p_0/4$.
Both choices precede the choice of $t$. Put $u=\log(t/K)$ and use the
stopped event
\begin{equation}
 F_u=\{\tau_u<\infty,\ S_j>0\ (1\le j\le\tau_u),\
           V_{\tau_u}\le K/2,\ O_u\le z\}.
 \label{eq:stopped-event}
\end{equation}
This event belongs to $\mathcal F_{\tau_u}$. The survival event in
\eqref{eq:survival}, minus the two exceptional events just bounded,
is contained in $F_u$. The union bound thus gives $\Q(F_u)\ge p_0/2$.

For every initial point with $2K\le|y|\le4K$, the pathwise inequality
\eqref{eq:radial-lower} yields on $F_u$
\begin{align}
 |X_j|&\ge\frac{3K}{2}e^{S_j}>K\qquad(1\le j\le\tau_u),
 \nonumber\\
 |X_{\tau_u}|&>\frac{3K}{2}\frac tK>t.\label{eq:excursion-path}
\end{align}
Thus $F_u$ forces a visit to $D_t$ before $C$. Since
$S_{\tau_u}\le u+z$ on $F_u$, the stopped identity
\eqref{eq:stopped-change} gives
\begin{equation}
 \Pp(F_u)\ge e^{-\alpha(u+z)}\Q(F_u)
 \ge\frac{p_0}{2}e^{-\alpha z}(K/t)^\alpha.
 \label{eq:excursion-likelihood}
\end{equation}
This proves \eqref{eq:annulus-bound}. The initial annulus lies outside
both target sets. Under $\Pp$, $\tau_u$ need not be almost surely finite;
\eqref{eq:stopped-change} already restricts to the finite-crossing event.
\end{proof}

\subsection{Stationary last-exit transfer}

\begin{lemma}[Last-exit inequality]\label{lem:last-exit}
For disjoint Borel sets $C,D$ and
$\tau_C^+=\inf\{n\ge1:X_n\in C\}$, a stationary chain satisfies
\begin{equation}
 \pi(D)\ge\int_C\pi(\dd x)\,\Pp_x(\tau_D<\tau_C^+).
 \label{eq:last-exit-main}
\end{equation}
\end{lemma}

\begin{proof}
Fix $N\ge1$. Among trajectories with $X_N\in D$, classify those having
a visit to $C$ at a time $0,\ldots,N-1$ according to their last such
visit. These events are disjoint. If the last visit occurs at $N-j$,
stationarity and the Markov property express its probability as
$\int_C\pi(\dd x)\Pp_x(X_j\in D,\ j<\tau_C^+)$. Consequently
\begin{align}
 \pi(D)
 &\ge\int_C\pi(\dd x)\sum_{j=1}^N
               \Pp_x(X_j\in D,\ j<\tau_C^+)\nonumber\\
 &\xrightarrow[N\to\infty]{}
 \int_C\pi(\dd x)\E_x
     \left[\sum_{j\ge1}\ind_{\{X_j\in D,\ j<\tau_C^+\}}\right]
 \nonumber\\
 &\ge\int_C\pi(\dd x)\Pp_x(\tau_D<\tau_C^+).
 \label{eq:last-exit-proof}
\end{align}
The passage to the infinite sum uses monotone convergence. The stationary
integral is unnormalized; no return-time or regeneration factor is omitted.
\end{proof}

\begin{proposition}[Strictly positive tail constant]\label{prop:positive}
For the constant from \cref{prop:tail-exists},
$C_{\eta,\theta}>0$.
\end{proposition}

\begin{proof}
Take $K,z$ from \cref{lem:excursion}, $C=[-K,K]$, and
$J=\{y:2K\le|y|\le4K\}$. Lemma~\ref{lem:minorization}, with
$R=K$ and $T=4K$, gives
\begin{equation}
 \inf_{x\in C}P_\theta(x,J)
 \ge d_K:=4K\varepsilon_{K,4K}>0.\label{eq:compact-annulus}
\end{equation}
Here $4K$ is the Lebesgue measure of the two-component annulus $J$.
The Markov property after the first step and \eqref{eq:annulus-bound}
yield, uniformly for $x\in C$,
\begin{equation}
 \Pp_x(\tau_{D_t}<\tau_C^+)
 \ge d_K\frac{p_0}{2}e^{-\alpha z}K^\alpha t^{-\alpha},
 \qquad t>4K.\label{eq:central-excursion}
\end{equation}
Full support implies $\pi(C)>0$. Applying \cref{lem:last-exit} now gives
\begin{equation}
 \pi(D_t)\ge c_{\eta,\theta}t^{-\alpha},\qquad
 c_{\eta,\theta}:=\pi(C)d_K\frac{p_0}{2}e^{-\alpha z}K^\alpha>0,
 \quad t>4K.\label{eq:positive-lower-limit}
\end{equation}
Thus $\liminf_{t\to\infty}t^\alpha\pi(D_t)\ge c_{\eta,\theta}>0$.
Combining this with the existing finite limit in
\cref{prop:tail-exists} proves
$C_{\eta,\theta}\ge c_{\eta,\theta}>0$.
\end{proof}

\section{Exact moment threshold and estimation error}\label{sec:moments}

For a nonnegative random variable $R$ and $p>0$, Tonelli's theorem gives
\begin{equation}
 \E R^p=p\int_0^\infty t^{p-1}\Pp(R>t)\dd t.
 \label{eq:tail-moment}
\end{equation}
The positive limit in \eqref{eq:tail-main} makes this integral diverge
for every $p\ge\alpha$, including the logarithmic divergence at
$p=\alpha$. Proposition~\ref{prop:invariant} proves finiteness for
$0<p<\alpha$.

A bounded translation preserves the exact tail constant. For $t>|\theta|$,
the triangle inequality gives
\begin{equation}
 \Pp(|X|>t+|\theta|)
 \le\Pp(|X-\theta|>t)
 \le\Pp(|X|>t-|\theta|).\label{eq:shift-sandwich}
\end{equation}
Multiply by $t^\alpha$ and use
$t^\alpha/(t\pm|\theta|)^\alpha\to1$. Both endpoints converge to
$C_{\eta,\theta}$, so the estimation-error tail has the same limit.
The same tail integral, together with bounded-shift moment inequalities,
gives its identical moment threshold. Since $\alpha<2$, both second
moments are infinite. This completes the proof of \cref{thm:main}.

\section{The concrete nonzero-signal example}\label{sec:signal}

At $\theta=1$, the labels are $Y=A+\delta\tanh A+E$, with
\begin{equation}
 \E[AY]=1+\delta\E[A\tanh A]>1.\label{eq:nonzero-signal}
\end{equation}
They contain a nonconstant input-dependent regression signal.
More generally, the population risk has the unique minimizer $\theta$:
\begin{align}
 F_\theta(x):=\E f_\theta(x;A,E)
 &=\frac12+\frac12\E[\sigma(Ax)-\sigma(A\theta)]^2\nonumber\\
 &\ge\frac12+\frac12(x-\theta)^2.\label{eq:risk}
\end{align}
Here independence and the mean-zero noise remove the cross term, and
$\sigma'\ge1$ together with $\E A^2=1$ proves the inequality.
At $x=\theta$, the risk is exactly $1/2$, proving uniqueness.

At $\eta=1$, the arithmetic in \eqref{eq:q-proof} sharpens to
\begin{equation}
 \E\Lip(\Psi_\theta)<\frac{97}{100}+\frac{10007}{10^6}
 =0.980007,
 \quad h_1(1)=\frac4{\sqrt{2\pi e}}<0.97,
 \quad h_1(2)=2.\label{eq:eta-one}
\end{equation}
In particular, $\theta=\eta=1$ is a genuine nonzero-teacher example
with an exact positive stationary estimation-error tail of index
$\alpha(1)\in(1,2)$, uniquely specified by $\E|1-A^2|^{\alpha(1)}=1$.

\section{The precise symmetry statement}\label{sec:symmetry}

Oddness of $\sigma$ and evenness of $\sigma'$ imply the exact identity
\begin{equation}
 \Psi_{-\theta,A,-E}(-x)=-\Psi_{\theta,A,E}(x).
 \label{eq:reflection}
\end{equation}
The additive noise is symmetric and the invariant law is unique, so
$\pi_{\eta,-\theta}$ is the reflection of $\pi_{\eta,\theta}$.
At $\theta=0$, one law is therefore symmetric, and for every $t>0$,
$\pi_{\eta,0}((t,\infty))=\pi_{\eta,0}((-\infty,-t))
=\pi_{\eta,0}(|x|>t)/2$. Each one-sided tail constant is
$C_{\eta,0}/2$. Reflection also gives
$C_{\eta,-\theta}=C_{\eta,\theta}$ for the absolute tails.

At fixed $\theta\ne0$, reflection changes the parameter and is not a
symmetry of that model. In fact exact stationary symmetry about zero is
impossible. Stationarity, the finite first moment, and averaging out the
fresh mean-zero noise give
\begin{equation}
 0=\E\!\left[A\sigma'(AX)
                   \bigl(\sigma(AX)-\sigma(A\theta)\bigr)\right].
 \label{eq:mean-stationarity}
\end{equation}
Both terms are integrable because $\sigma'\le1+\delta$,
$|\sigma(v)|\le(1+\delta)|v|$, and $\E_\pi|X|<\infty$.
If $\pi$ were symmetric, $\E[A\sigma'(AX)\sigma(AX)]=0$ since its
integrand is odd in $X$ for each fixed $A$. But
$A\sigma'(AX)\sigma(A\theta)$ has the strict sign of $\theta$ almost
surely when $A\ne0$, contradicting \eqref{eq:mean-stationarity}.
Equality of asymptotic one-sided constants would require a separate
signed-tail argument; it is not deduced here.

\section{Attribution, source pins, and scope}\label{sec:attribution}

The original Math Brain T228 materials already contain scalar smooth
asymptotically linear results, a leaky-ReLU analysis, and finite-dimensional
radial-tail arguments under their stated stability assumptions
\cite{BrainT1,BrainT2,BrainT3Setup,BrainT3Stability,BrainT3Tail,BrainT3Positive}.
The forward tilted-walk, overshoot, and stationary last-exit mechanism in
\cref{sec:positive} is credited specifically to
\texttt{claude-b1e26540}'s \texttt{T3\_STEP3.md}. This manuscript retains
that attribution and rechecks the changed teacher-dependent label,
remainder, complete-step tilt, and conditional Gaussian mean. The direct
exponential-supermartingale survival bound and explicit finite-horizon
last-exit sum are supplied in the campaign's nonzero-teacher proof.
No discovery credit is reassigned.

The principal campaign source is
\texttt{T228\_NONZERO\_TEACHER.md}, version 1, 23 October 2026, by the
campaign's A09 extension and verification. The baseline is
\texttt{T228\_EXPLICIT\_CLOSURE.md}, version 2 of the same date. The
baseline supplies an explicit average-contraction regime in addition to
the prior stability-based arguments. The internal campaign proof-source
SHA-256 values are recorded in \cref{app:provenance}. Independent internal
audits are research checks; they are not external peer review or a
journal acceptance.

The sole external exact-tail theorem used in this proof is
\cite[Theorem 3.1]{JO2012Thm31}, in the primary arXiv v5 PDF. The arXiv
HTML rendering numbers the corresponding Section 3 statement differently;
the citation here is to the printed PDF theorem. The background
SGD reference \cite{Gurbuzbalaban2021} motivates the problem and does not
replace the hypothesis verification above.

The result is confined to the concrete scalar model and parameter
rectangle in \eqref{eq:model-parameters}. It is not a theorem for arbitrary
step sizes, all bounded activations, arbitrary nonlinear networks, or
deep-network architectures. It does not establish monotonicity in
$\eta/b$, vary the batch size, or certify first-in-the-literature priority.
No simulation, numerical tail fit, or quadrature substitutes for any part
of the proof. Unreviewed campaign extensions are excluded from this
manuscript.

\appendix
\section{Exact arithmetic and proof-source identification}\label{app:provenance}

The following identities are rational arithmetic checks, not numerical
approximations used in place of proof.
\begin{center}
\small
\begin{tabularx}{\textwidth}{@{}X>{\raggedleft\arraybackslash}p{0.37\textwidth}@{}}
\toprule
Quantity & Exact value \\
\midrule
Conservative mixed-label bound $3+2\delta$ & $1501/500=3.002$ \\
Expected derivative perturbation $4\delta+3\delta^2+2\delta(3+2\delta)$
 & $10007/10^6=0.010007$ \\
Uniform contraction upper bound $q$ & $99010707/10^8=0.99010707$ \\
Contraction margin $1-q$ & $989293/10^8=0.00989293$ \\
Upper bound at $\eta=1$ & $980007/10^6=0.980007$ \\
Minimum of $h_\eta(2)$ on the interval & $19603/10000=1.9603$ \\
Squared rational normal-bound margin in \eqref{eq:rational-normal}
 & $7532467/125000000>0$ \\
\bottomrule
\end{tabularx}
\end{center}

\paragraph{Campaign proof inputs.}
The exact source content used for this typesetting is identified by
SHA-256; these are local proof revisions, not claims of publication.
\noindent\emph{Main proof: \texttt{T228\_NONZERO\_TEACHER.md}, version 1.}\\
\sourcehash{f95fa6fade82ba0dca4dab8e311fc61bb57b3267f3953b465c94228cb1315933}

\noindent\emph{Baseline: \texttt{T228\_EXPLICIT\_CLOSURE.md}, version 2.}\\
\sourcehash{a1dbdbb1e99a85749143f44dd8b7d7e264fb550d285c5667f69f1a35ff87c38f}

\paragraph{Math Brain content pins.}
The references below give the source paths and immutable Git blob SHA-1
pins recovered in \texttt{a01/SOURCE\_MANIFEST.json}. A branch URL is a
convenience link; the accompanying blob pin identifies the exact retrieved
content. The repository is
\href{https://github.com/<private-repo>}
{\texttt{<private-repo>}}.
The path prefix for the T228 entries is
\texttt{06\_math\_research/frontier/rows/T228/}.
The dossier itself has Git blob pin
\sourcehash{5d95a5f8a9d830ca673e75321d66f66d9dccad45}.

% Reserved appendix hook: reviewed campaign extensions may be added here
% only after the coordinator explicitly supplies a frozen, audited source.
% No A14/A15 theorem is imported in this version.

\begin{thebibliography}{10}
\raggedright
\addcontentsline{toc}{section}{References}

\bibitem{JO2012Thm31}
P.~R. Jelenkovi\'c and M.~Olvera-Cravioto.
\emph{Implicit Renewal Theory and Power Tails on Trees}.
arXiv:1006.3295v5, 31 May 2012.
Theorem 3.1, printed p.~5.
\href{https://arxiv.org/abs/1006.3295v5}{Versioned primary record};
\href{https://arxiv.org/pdf/1006.3295}{primary PDF}.

\bibitem{Gurbuzbalaban2021}
M.~G\"urb\"uzbalaban, U.~\c{S}im\c{s}ekli, and L.~Zhu.
The Heavy-Tail Phenomenon in SGD.
In \emph{Proceedings of the 38th International Conference on Machine
Learning}, PMLR \textbf{139}, 3964--3975, 2021.
\href{https://proceedings.mlr.press/v139/gurbuzbalaban21a.html}
{Primary proceedings record};
\href{https://arxiv.org/abs/2006.04740}{arXiv:2006.04740}.

\bibitem{BrainT1}
\texttt{codex-w688} (Math Brain).
\emph{T228 target 1: scalar asymptotically linear SGD}, version 1.
Path: \texttt{w688\_target1/proofs/T228\_A\_v1.md}.
Git blob pin: \sourcehash{6056b253ff1d4b9eb2411d6dec7d25c5ce8d7523}.
\href{https://github.com/<private-repo>/blob/main/06_math_research/frontier/rows/T228/w688_target1/proofs/T228_A_v1.md}
{Source path}.

\bibitem{BrainT2}
\texttt{claude-4af6ff1a} (Math Brain).
\emph{T228 target 2: scalar leaky-ReLU SGD has an exact power tail},
version 2, 21 October 2026.
Path: \texttt{claude\_4af6ff1a/T228\_B\_v2.md}.
Git blob pin: \sourcehash{b7b358985ff34a9ea3fa7b35339047f0fe793822}.
\href{https://github.com/<private-repo>/blob/main/06_math_research/frontier/rows/T228/claude_4af6ff1a/T228_B_v2.md}
{Source path}.

\bibitem{BrainT3Setup}
\texttt{claude-b1e26540} (Math Brain).
T228 target 3 setup, 22 October 2026.
Path: \texttt{claude\_b1e26540/T3\_SETUP.md}.
Git blob pin: \sourcehash{de30be030e66bf30700ba818a5155c9d6c4e3562}.
\href{https://github.com/<private-repo>/blob/main/06_math_research/frontier/rows/T228/claude_b1e26540/T3_SETUP.md}
{Source path}.

\bibitem{BrainT3Stability}
\texttt{claude-b1e26540} (Math Brain).
T228 target 3, step 1, 22 October 2026.
Path: \texttt{claude\_b1e26540/T3\_STEP1.md}.
Git blob pin: \sourcehash{c169b975cd3a5d8ed1f554214d010da1c587e96d}.
\href{https://github.com/<private-repo>/blob/main/06_math_research/frontier/rows/T228/claude_b1e26540/T3_STEP1.md}
{Source path}.

\bibitem{BrainT3Tail}
\texttt{claude-b1e26540} (Math Brain).
T228 target 3, step 2, 22 October 2026.
Path: \texttt{claude\_b1e26540/T3\_STEP2.md}.
Git blob pin: \sourcehash{04d08036877af25a6d083247fdbfbf834ae0206e}.
\href{https://github.com/<private-repo>/blob/main/06_math_research/frontier/rows/T228/claude_b1e26540/T3_STEP2.md}
{Source path}.

\bibitem{BrainT3Positive}
\texttt{claude-b1e26540} (Math Brain).
\emph{T228 target 3, step 3: the tail constant is strictly positive},
22 October 2026.
Path: \texttt{claude\_b1e26540/T3\_STEP3.md}.
Git blob pin: \sourcehash{a3adfc7537b6a8f51b796539c5063b07a23fca39}.
\href{https://github.com/<private-repo>/blob/main/06_math_research/frontier/rows/T228/claude_b1e26540/T3_STEP3.md}
{Source path}.

\end{thebibliography}
\end{document}
